\chapter{Discussão}\label{cap:discussao}

\section{Resposta às questões de pesquisa}

\textbf{Q1 --- Como estruturar um servidor MCP que cubra o fluxo completo de
PCB no KiCad?} A arquitetura em duas camadas, apresentada no
Capítulo~\ref{cap:arquitetura}, mostrou-se adequada: a camada TypeScript isola
o protocolo e permite testes sem o KiCad; a camada Python concentra o domínio
e a comunicação com os três mecanismos de automação (SWIG, IPC e
\texttt{kicad-cli}). O catálogo resultante cobre da criação do projeto à
exportação de fabricação em 24 módulos (Capítulo~\ref{cap:catalogo}).

\textbf{Q2 --- Como organizar e tornar descobríveis centenas de operações?} A
combinação de registro por categorias, lista de essenciais e busca por
palavra-chave cobre 74,2\,\% das ferramentas. O restante permanece como dívida
de descoberta controlada por teste. A experiência do projeto indica que a
descoberta não é um detalhe: sem ela, um catálogo de 233 operações torna-se
difícil de usar por agentes.

\textbf{Q3 --- A implementação resiste a falhas previsíveis de integração?} Os
mecanismos de correlação de requisições, descarte de respostas tardias,
tempos-limite por comando e sessão fixada por \emph{backend} atacam
exatamente as classes de falha observadas no histórico do projeto. As suítes
registraram 2.412 de 2.451 casos aprovados na execução local; as 8 falhas
remanescentes estão associadas ao ambiente de validação e a caminhos legados,
conforme análise do Capítulo~\ref{cap:resultados}.

\textbf{Q4 --- Que evidências de qualidade e uso real podem ser obtidas?} O
estudo de caso produziu um conjunto concreto de evidências: caracterização do
projeto (97 \emph{footprints}, 83 redes, 9 folhas), verificação de regras com
11 violações residuais e zero itens não conectados, renderizações 3D e
artefatos de fabricação. A verificação, em especial, revelou pendências reais
--- conexões térmicas insuficientes e anéis anulares críticos --- que
precisam de correção antes da fabricação.

\section{Comparação com a literatura}

Os achados dialogam com três resultados da literatura. Primeiro, a
importância da interface agente--computador \citep{yang2024sweagent} foi
confirmada na prática: decisões como expor uma ferramenta por operação, manter
uma lista de essenciais e devolver erros estruturados decorrem diretamente da
necessidade de tornar as ações compreensíveis para o agente. Segundo, a
complementaridade em relação à adaptação de modelos proposta por
Liu et al.~\cite{liu2023chipnemo}: o servidor não altera o modelo e, por isso, beneficia
qualquer assistente compatível; os dois caminhos podem ser combinados. Por
fim, o trabalho explora o contrato padronizado do MCP \citep{mcp2025spec} em
um domínio de engenharia com operações destrutivas, ampliando os relatos
predominantes de integração em domínios de software.

\section{Contribuições}

As contribuições do trabalho são:

\begin{enumerate}[leftmargin=1.4em,itemsep=0.2em]
  \item a análise documentada do \texttt{mcp-kicad-vscode}~\citep{kicadmcpserver},
        servidor MCP de código aberto com cobertura ampla do fluxo de PCB no
        KiCad (233 ferramentas em 24 módulos), incluindo integrações de
        fabricação e cadeia de suprimentos;
  \item a caracterização de sua arquitetura em duas camadas, com protocolo
        interno correlacionado e abstração de \emph{backends} (SWIG/IPC/CLI);
  \item a avaliação do modelo de descoberta em três mecanismos (categorias,
        essenciais e busca), com salvaguarda automatizada contra regressão de
        cobertura;
  \item a consolidação de uma base de verificação com 2.451 casos de teste e
        integração contínua multiplataforma, descrita como referência para
        integrações semelhantes;
  \item um estudo de caso completo, conduzido pelo autor, com métricas
        verificáveis de um projeto de quatro camadas, incluindo a análise
        crítica das violações residuais de regras de projeto.
\end{enumerate}

\section{Limitações}\label{sec:limitacoes}

\begin{enumerate}[leftmargin=1.4em,itemsep=0.25em]
  \item \textbf{Integração IPC não exercitada.} No ambiente de validação, a
        API IPC do KiCad estava desabilitada na configuração do usuário e sem
        soquete ativo; os resultados refletem o caminho SWIG e o
        \texttt{kicad-cli}. A validação do fluxo em tempo real permanece como
        o principal trabalho futuro.
  \item \textbf{Cobertura de descoberta parcial.} Sessenta ferramentas não
        são localizáveis por busca; embora a dívida esteja congelada por
        teste, ela limita a autonomia do agente.
  \item \textbf{Falhas de teste sensíveis ao ambiente.} Oito casos falharam
        na execução local, todos dependentes de subprocessos
        \texttt{kicad-cli} ou de caminhos legados; o comportamento em
        ambientes de CI limpos pode divergir.
  \item \textbf{Pendências de fabricação no caso estudado.} As 11 violações
        residuais --- 9 erros, incluindo anel anular crítico --- impedem a
        leitura do projeto como pronto para fabricação; são, ao mesmo tempo,
        evidência do valor da verificação.
  \item \textbf{Avaliação restrita.} Não houve estudo com usuários, medição
        de tempo de projeto ou comparação controlada com fluxo manual ou
        roteirizado; as métricas são descritivas e de um único projeto.
  \item \textbf{Escopo do modelo.} O trabalho não avalia a qualidade de
        raciocínio do modelo de linguagem; assume que o cliente MCP é capaz
        de planejar e invocar ferramentas.
\end{enumerate}

\section{Ameaças à validade}

\textbf{Validade de construto}: contagens de ferramentas e de testes são
proxis de capacidade e confiabilidade; a utilidade percebida não foi medida.
\textbf{Validade interna}: parte das falhas de teste é atribuível ao ambiente
\emph{snap} do KiCad na estação de validação, o que recomenda replicação em
ambiente de CI padrão. Além disso, os dados do estudo de caso refletem a
revisão vigente do projeto na data da verificação; projetos evoluem.
\textbf{Validade externa}: a generalização para outros \emph{softwares} de
EDA exige esforço de porte; contudo, o desenho de duas camadas com
\emph{backend} abstrato é replicável. \textbf{Confiabilidade}: os
procedimentos são reprodutíveis (comandos e logs preservados); as suítes são
determinísticas, exceto pelos casos sensíveis a ambiente.

\section{Implicações práticas}

Três implicações merecem destaque. Na \textbf{educação em engenharia}, o
servidor permite que estudantes conduzam verificações de regras e
entendam o fluxo de fabricação com assistência; na \textbf{prototipagem
industrial}, o encadeamento projeto--verificação--fabricação reduz passos
manuais e padroniza a troca de informação com fornecedores; na
\textbf{qualidade de projetos}, a combinação de DRC/ERC e artefatos de
fabricação automatizados cria trilhas de auditoria verificáveis.

\section{Trabalhos futuros}

\begin{enumerate}[leftmargin=1.4em,itemsep=0.2em]
  \item habilitar e validar o fluxo IPC em tempo real, com testes de
        integração ponta a ponta na interface gráfica;
  \item indexar as 60 ferramentas restantes e avaliar busca semântica de
        operações;
  \item implementar camada de políticas de permissão e auditoria para
        operações destrutivas;
  \item reduzir a sensibilidade ambiental da suíte e ampliar a cobertura de
        formatos de arquivo legados;
  \item conduzir estudos com usuários para medir tempo de projeto, taxa de
        retrabalho e percepção de utilidade;
  \item avaliar a solução com modelos e clientes MCP distintos, incluindo
        execução local de modelos;
  \item evoluir o caso de estudo para fechamento das violações residuais e
        fabricação completa da placa.
\end{enumerate}

\section{Síntese do capítulo}

A discussão conectou os resultados às questões de pesquisa e à literatura,
explicitou as contribuições, delimitou limitações e ameaças à validade e
propôs uma agenda de continuidade. O capítulo final retoma o percurso e
encerra o trabalho.
